\eabstract{
    \par The phenomenon of quantum entanglement carries profound physical significance, and testing entanglement between particles under high-energy collider conditions can provide a novel perspective for understanding its intrinsic nature. Quantum tomography, by establishing a functional relationship between differential scattering cross-sections and spin correlation matrices, have offered a practical computational tool for experimentally determining the presence of entanglement. With the rise of artificial intelligence and the widespread application of machine learning methods, testing entanglement in particle pairs within particle physics experiments has gradually become feasible.

    \par This study primarily employs the OmniLearn generative model, constructed based on diffusion models and Transformer architectures, to reconstruct the kinematic information of final-state particles in Monte Carlo simulations of $pp \to \tau^+ \tau^- \to (\pi^+ \bar{\nu}_\tau)(\pi^- \nu_\tau)$. After performing reference frame transformations and statistical analysis of all events, the angular distributions between the momentum direction of the $pi$ mesons and the helicity basis are obtained, enabling the calculation of all matrix elements of the two-qubit ensemble spin correlation matrix for $\tau^+ \tau^-$. Utilizing the entanglement criterion derived from the Peres-Horodecki criterion, $\mathcal{O}_E^\pm=0.1688>0$, it is confirmed that the spins of the $\tau^+ \tau^-$ pair in this process exhibit entanglement.

    \par The machine learning method employed in this research generates kinematic information for each particle that aligns well with the reconstructed-level data from Monte Carlo simulations, demonstrating significant advantages over traditional missing-mass calculator approaches. This advantage serves as strong evidence for the reliability of the spin correlation matrix calculations and the entanglement determination presented in this work.

    }{Quantum Entanglement; Collider; $\tau^+ \tau^-$; Generative Model}
